\begin{abstract}
Digital-twin network planning needs path-loss models for areas without
measurements. Tile-based methods fit a log-distance exponent and intercept
per map tile and transfer them to unmeasured tiles by descriptor similarity.
We study this transfer across cities on ray-traced corpora:
\nCities{} geo-referenced cities with \nSites{} rooftop sites at
\SI{\freqGHz}{\giga\hertz} (\nSamples{} links) and
\SI{\freqTwoGHz}{\giga\hertz}, a random-siting replicate and independently
modelled external scenes. The tile labels are imprecise: their
cluster-robust uncertainty is \seRatio{} times the naive one, and they change
with the serving sites and the dynamic range of the data. Learned regressors
transfer them better than the source median in most cities, though not
significantly. Yet even perfect tile labels give a path-loss error of
\SI{\netOracle}{\decibel} and the correct serving site for only
\assocOracle\% of the pixels; composing the parameters of the tiles a link
crosses helps with exact labels, not once they are transferred. A site-aware hurdle model that predicts each
link, and whether it exists, from its obstruction and diffraction geometry
reduces the error on unseen cities to \SI{\netLink}{\decibel}, raises
serving-site accuracy to \assocLink\% and cuts the coverage lost when
planning five sites from \regretGbdt\% to \regretLink\%, in every held-out
city. On external scenes with other materials, a survey of 25 random tiles
calibrates it to the accuracy of exact tile labels. A city-level
certificate lets it promise \certClaimLink\% of the pixels as covered at a
guaranteed risk, against \certClaimGbdt\% for tile transfer.
All results are reproducible by one command.
\end{abstract}

\begin{keyword}
digital twin \sep network planning \sep path-loss modelling \sep transfer
learning \sep ray tracing \sep coverage certification \sep empirical Bayes
\end{keyword}
